\documentclass[a4paper,leqno,12pt]{amsart}
%
\usepackage{amsmath,amsthm}
\usepackage{graphicx}
\usepackage[margin=0.7in]{geometry}

\usepackage{enumitem}

\usepackage{comment}
\usepackage{multirow}
\usepackage{tikz}
\usepackage{wrapfig}
\usepackage[autostyle=true]{csquotes}

\usepackage{fancyhdr}
\pagestyle{fancy}
\fancyhf{} % clear all header/footer fields
\fancyfoot[C]{\bf ************************************} % 

\newlist{Case}{enumerate}{1}
\setlist[Case]{label=\bf{Case ~\arabic*:}}
%
\begin{document}
\begin{center}
{{\textbf{ AZIM PREMJI UNIVERSITY, BHOPAL}\\
First Semester  2026-2027\\
Tutorial Sheet-3}}
\end{center}

\noindent Course No. MTH-113 \hfill  Course title: Introduction to Quantitative Thinking
\vspace{.3cm}
\hrule 
\vspace{.3cm}


\begin{enumerate}
\setlength\itemsep{1em}
\item What will be the unit digit of the squares of the following numbers?\\
(i) $81 \quad$ (ii) $272 \quad$ (iii) $799 \quad $ (iv) $3853 \quad $ (v) $1234 \quad $ (vi) $26387$

\item The squares of which of the following would be odd numbers? \\
(i) $431 \quad$ (ii) $2826 \quad$ (iii) $7779 \quad $ (iv) $82004$

\item Observe the following pattern and find the missing digits.
\begin{eqnarray*}
 11^2 &=& 121 \cr
 101^2 &=& 10201 \cr 
 1001^2 &=& 1002001 \cr
 100001^2 &=& 1 ......... 2 ......... 1 \cr
 10000001^2 &=& ...........................
\end{eqnarray*}

\item Using the given pattern, find the missing numbers.\\
$1^2 + 2^2 + 2^2 = 3^2$\\
$3^2 + 4^2 + 12^2 = 13^2$\\
$4^2 + 5^2 + \_^2 = 21^2$\\
$5^2 + \_^2 + 30^2 = 31^2$\\
$6^2 + 7^2 + \_^2 = \_^2$

\item How many numbers lie between squares of the following numbers?\\
(i) $12 \ {\rm and} \ 13 \quad$ (ii) $25 \ {\rm and} \ 26 \quad$ (iii) $99 \ {\rm and} \ 100$


\item Verify that the following equality holds for any integer $m \geq 0$:
\[
(2m)^2 + (m^2 - 1)^2 = (m^2 + 1)^2
\]
In this case, the triplet $(2m, m^2 -1, m^2 +1)$ is known as $\textit{Pythagorean triplet}$.

\item Write a Pythagorean triplet whose smallest member is 8.

\item Write a Pythagorean triplet whose one member is:\\
(i) $6 \quad $ (ii) $14 \quad$ (iii) $16 \quad$ (iv) $18$ 

\item Find the smallest square number which is divisible by each of the numbers 6, 9 and 15. 

\item What could be the possible ‘one’s’ digits of the square root of each of the following
numbers? \\
(i) $9801 \quad$ (ii) $99856 \quad$ (iii) $998001$

\item Without doing any calculation, find the numbers which are surely not perfect squares.\\
(i) $153 \quad$ (ii) $257 \quad$ (iii) $408 \quad$ (iv) $441$

\item Find the square roots of the following numbers by the Prime Factorisation Method.\\
(i) $153 \quad$ (ii) $1764 \quad$ (iii) $4096 \quad$ (iv) $7744$

\item For each of the following numbers, find the smallest whole number by which it should
be multiplied so as to get a perfect square number. Also find the square root of the
square number so obtained.\\
(i) $252 \quad$ (ii) $1008 \quad$ (iii) $2028 \quad$ (iv) $1458$

\item Find the square root of each of the following numbers by Division method.\\
(i) $2304 \quad$ (ii) $4489 \quad$ (iii) $3481 \quad$ (iv) $5776$

\item Find the least number which must be subtracted from each of the following numbers
so as to get a perfect square. Also find the square root of the perfect square so
obtained.\\
(i) $402 \quad$ (ii) $1989 \quad$ (iii) $3250 \quad$ (iv) $825$

\item Find the least number which must be added to each of the following numbers so as
to get a perfect square. Also find the square root of the perfect square so obtained.\\
(i) $525 \quad$ (ii) $1750 \quad$ (iii) $252 \quad$ (iv) $1825$

\item Find the smallest number by which each of the following numbers must be multiplied
to obtain a perfect cube.\\
(i) $243 \quad$ (ii) $256 \quad$ (iii) $72 \quad$ (iv) $675$

\item You are told that 1,331 is a perfect cube. Can you guess without factorisation what
is its cube root? Similarly, guess the cube roots of 4913, 12167, 32768.

\item Find the value of.\\
(i) $(3^0  + 4^{-1}) \times 2^2 \quad $ (ii) $(2^{-1} \times 4^{-1}) \div 2^{-2} \quad$ (iii) $\left(\frac{1}{2} \right)^{-2} + \left(\frac{1}{3} \right)^{-2} + \left(\frac{1}{4} \right)^{-2} \quad$ (iv) $(3^{-1} + 4^{-1} + 5^{-1})^0$\\
(v) $\frac{8^{-1} + 5^3}{2^{-4}} \quad$ (vi) $\left\{\left(\frac{1}{3}\right)^{-1} -\left(\frac{1}{4}\right)^{-1} \right\}^{-1}$

\item Express the following numbers in usual form.\\
(i) $3.02 \times 10^{-6} \quad$ (ii) $4.5 \times 10^4 \quad$ (iii) $3 \times 10^{-8} \quad $ (iv) $3.61492 \times 10^6$

\item Express the number appearing in the following statements in standard form.\\
\begin{enumerate}
 \item 1 micron is equal to $\frac{1}{1000000}$ m
 \item Charge of an electron is 0.000,000,000,000,000,000,16 coulomb.
 \item Size of a bacteria is 0.0000005 m
 \item Size of a plant cell is 0.00001275 m
 \item Thickness of a thick paper is 0.07 mm
\end{enumerate}



\begin{comment}
For each $n = 2$, the group $\pi_1(CP^2 - C)$ has cohomological dimension 2 (for instance, the complement $CP^2 - C$ is known to be a $K(\pi_1,1)$-space, and has the homotopy type of a 2-dimensional complex). In particular, when $n  = 2$, the group $\pi_1(CP^2 - C)$ is the quotient of the 4-strand Artin pure braid group by its center (isomorphic to $F_3$ o $F_2$). In this $n = 2$ instance, the group G presented in Definition 4.3 is
\[
G = \{g_0,g_1,g_0' ,g_1' | [g_0,g_1], [g_0' ,g_1' ], [g_0,g_0' ], [g_1,g_1' ]\}
\]
a (right-angled Artin) group of cohomological dimension 2. From the preceding Proposition 4.1, it appears
that the action of the free group $F = <g_0'',g_1''>$ on (the generators of) G is by conjugation, so is trivial in 201. This would imply that $H_3(G \times F_2; Z) \neq 0$, hence $cd (G \times F_2) > 2$
\end{comment}


\end{enumerate}


%\vspace*{0.2cm}
%\begin{center}
%{\bf ************************************}
%\end{center}

\end{document}